% ============================================================
% Dịch vụ nhận dạng giọng nói (Voice Rephraze Service)
% Dùng trong: Chương 3 (Thiết kế) và Chương 4 (Cài đặt)
% ============================================================

% --- THIẾT KẾ (được \input từ analysis-design.tex) ---

\newcommand{\designVoice}{%

\section{Thiết kế dịch vụ nhận dạng giọng nói}

Dịch vụ Voice Rephraze đảm nhận hai chức năng chính: chuyển đổi giọng nói
thành văn bản (Speech-to-Text) và sinh mô tả sản phẩm. Cả hai được triển khai
trên FastAPI với AsyncOpenAI client, tận dụng khả năng xử lý bất đồng bộ để
gọi API OpenAI hiệu quả hơn.

Đối với module STT, dự án sử dụng OpenAI gpt-4o-transcribe --- mô hình cloud
API hỗ trợ nhận dạng tiếng Việt đa vùng miền với nhiều định dạng audio (WAV,
MP3, FLAC, OGG, WebM, M4A). Module sinh mô tả sử dụng GPT-4o-mini với
temperature 0.75 và giới hạn 600 tokens. System prompt được thiết kế để AI
đóng vai copywriter thương mại điện tử, có khả năng trích xuất thông tin sản
phẩm từ text lộn xộn của STT --- bỏ qua từ lấp, lặp lại và lỗi phát âm ---
rồi viết 3--5 câu ngắn gọn theo phong cách được chọn.

Hệ thống hỗ trợ nhiều phong cách viết mô tả, được lưu trữ trong SQLite và
có thể mở rộng thông qua API. \Cref{tab:writing-styles} liệt kê các phong cách
mặc định.

\begin{figure}[htbp]
\centering
\fbox{\parbox{0.85\linewidth}{\centering\vspace{2cm}
\textit{Thêm hình: Sơ đồ luồng xử lý dịch vụ Voice Rephraze --- Audio đầu vào
$\rightarrow$ gpt-4o-transcribe (STT) $\rightarrow$ raw text $\rightarrow$
GPT-4o-mini + prompt engineering + style $\rightarrow$ mô tả sản phẩm hoàn chỉnh.
Thể hiện cả luồng quản lý phong cách viết từ SQLite.}
\vspace{2cm}}}
\caption{Luồng xử lý dịch vụ nhận dạng giọng nói}
\label{fig:voice-rephraze-flow}
\end{figure}

\begin{table}[htbp]
\centering
\caption{Các phong cách viết mô tả sản phẩm}
\label{tab:writing-styles}
\begin{tabular}{>{\raggedright\arraybackslash}p{0.25\linewidth} >{\raggedright\arraybackslash}p{0.6\linewidth}}
\toprule
\rowcolor{headerblue} \textcolor{white}{\textbf{Phong cách}} & \textcolor{white}{\textbf{Mô tả}} \\
\midrule
\rowcolor{rowgray} Văn minh & Lịch sự, trang trọng, ngôn ngữ chuẩn mực \\
Chuyên nghiệp & Kỹ thuật, rõ ràng, tập trung thông số \\
\rowcolor{rowgray} Thân thiện & Ấm áp, gần gũi, dễ tiếp cận \\
Văn hóa người Tày & Yếu tố văn hóa dân tộc, hình ảnh đặc trưng \\
\bottomrule
\end{tabular}
\end{table}

}% end \designVoice

% --- CÀI ĐẶT (được \input từ implementation.tex) ---

\newcommand{\implVoice}{%

\section{Cài đặt dịch vụ nhận dạng giọng nói}

\subsection{Speech-to-Text}

\begin{itemize}
\item \textbf{Mô hình:} OpenAI gpt-4o-transcribe (cloud API).
\item \textbf{Định dạng hỗ trợ:} WAV, MP3, FLAC, OGG, WebM, M4A.
\item \textbf{Ngôn ngữ:} Tiếng Việt (đa vùng miền).
\item \textbf{API:} \texttt{POST /stt} -- upload file audio, trả về text.
\end{itemize}

\subsection{Sinh mô tả sản phẩm}

\begin{itemize}
\item \textbf{Mô hình:} GPT-4o-mini (temperature 0.75, max 600 tokens).
\item \textbf{API:} \texttt{POST /gen} -- nhận \texttt{product\_description} (text từ STT)
  và \texttt{style} (phong cách viết).
\item \textbf{Prompt engineering:} System prompt hướng dẫn AI đóng vai copywriter
  thương mại điện tử, trích xuất thông tin sản phẩm từ text lộn xộn của STT,
  bỏ qua từ lấp, lặp lại, lỗi phát âm; viết 3--5 câu ngắn gọn theo phong cách yêu cầu.
\end{itemize}

\subsection{Lưu trữ phong cách viết}

Các phong cách viết được lưu trong SQLite (\texttt{styles.db}) với schema đơn giản
(id, name, description). Người dùng có thể thêm phong cách mới qua
\texttt{POST /add-style}.

}% end \implVoice
